\subsection{A restriction on identity relative seeds for the harmonic defect}

\begin{statement}{D23.P15}{Proved}{A normal relative seed cannot separate the two order-48 evaluations by nilpotent factorization}
\end{statement}
Let $B=C_2\amalg C_3$, $F=\widehat F(x,y)$, and
$x=(sr)^{-2}$, $y=sxs$ as in the supplied conventions. Put
$R_4=(\mathbb Z/4)[\omega]$, $\omega^2+\omega+1=0$, and use
the following two marked surjections onto groups of order $48$:
\[
 q_K(x)=(0,\omega),\quad q_K(y)=(1,1),\qquad
 q_H(x)=(0,\omega),\quad q_H(y)=(1,\omega)
 \quad\text{in }R_4\rtimes\langle\omega\rangle.
\]
If $U\triangleleft B$ is an open normal subgroup contained in $F$
and $q_K(U)$ is nilpotent, then
\[
 U\subseteq F'F^3,\qquad q_H(U)\subseteq(R_4,+).
 \tag{D23.P15.1}
\]
Consequently both maps kill $R_{\rm nil}(U)$. In particular, an
actual gentle pair $(1,f)$ with
$f\in R_{\rm nil}(U)$ necessarily satisfies
\[
 \Psi_3(f)/2\equiv0\pmod2,
 \qquad \rho(f)\equiv0\pmod3,
 \qquad a_2(f)=0.
 \tag{D23.P15.2}
\]
Thus an identity relative construction which forces the Kummer
target to factor through the relative pronilpotent quotient cannot
produce $a_2=1$ and $\Psi_3=0$. This statement does not exclude a
nonidentity relative automorphism, a different construction, or an
actual gentle element with those character values.

\begin{proof}
Write $V=(R_4,+)$ and $G=V\rtimes C_3$. Since $U$ is normal in
$F$ and $q_K$ is onto, $N=q_K(U)$ is normal in $G$. If its image
in $C_3$ were nontrivial, $N$ would contain an element with linear
part $\omega$ or $\omega^2$. Its commutators with $V$ generate
$(1-\omega)V=V$, because $1-\omega$ is a unit in $R_4$.
Normality would give $V\subseteq N$, and then $N=G$. But $G$
is not nilpotent: its order-three action on $V$ is nontrivial,
whereas distinct Sylow subgroups commute in a finite nilpotent
group. It follows that $N\subseteq V$.

The linear part of $q_K(u)$ is $\omega^{\operatorname{exp}_x(u)}$.
Therefore every $u\in U$ has $x$ exponent divisible by $3$.
Conjugation by $s$ exchanges $x$ and $y$. Since $sUs=U$, the
$y$ exponent of every $u\in U$ is divisible by $3$ as well.
This proves $U\subseteq F'F^3$. The linear part of $q_H(u)$ is
$\omega^{\operatorname{exp}_x(u)+\operatorname{exp}_y(u)}$, so
it is one. This proves (D23.P15.1), and $q_H(U)$ is abelian.

The universal property of the nilpotent residual now gives
$q_K(R_{\rm nil}(U))=q_H(R_{\rm nil}(U))=1$.
The supplied A.89--A.90 identify the half-character modulo two
with the scalar translation of $q_K(f)$ reduced modulo two;
it is therefore zero. The supplied FP.34 harmonic formula is
\[
 1-(1-\omega)\operatorname{trans}(q_H(f))
       =\omega^{\rho(f)}(-1)^{a_2(f)}.
\]
Its left side is one. The six phases
$\omega^j(-1)^e$, $j\in\mathbb Z/3$, $e\in\mathbb Z/2$, are
distinct in $R_4^\times$, so $\rho\equiv0\pmod3$ and $a_2=0$.
This proves (D23.P15.2).
\end{proof}

\paragraph{Source and check scope.}
The two marked quotients and their character formulas are the
supplied \texttt{DELTA\_ON\_GT/sections/A.tex}, statements
A.89--A.90, and \texttt{agent\_five/five\_point\_results.tex},
FP.34. The relative lifting theorem itself is supplied
\texttt{LM7\_PUSH/E\_relative\_cm\_lift\_agent.tex}, E.58;
its missing D.50 proof is not audited here. The argument above
is an independent proof of this restriction on its identity
relative application, with no assertion about other seeds.
